% Replace the sections ``Comparable systems'' and ``Synthesis of the state of
% the art'' in the State of the Art chapter with the following blocks.

\section{Comparable systems}

Several families of systems address parts of the problem considered in this
thesis, although few combine all of its requirements: heterogeneous French
source material, traceable collection, controlled metadata, and exploratory
analysis of territorial initiatives.

A first family consists of documentary catalogues and digital-library systems.
They organise resources through bibliographic or descriptive metadata and
provide faceted browsing by author, date, subject, institution, or place. Their
main strength is the use of stable metadata schemas and controlled vocabularies,
which improve consistency and make filtering possible. As noted in work on
taxonomies and controlled vocabularies, this standardisation reduces ambiguity
and supports reliable indexing \citep{hedden2016taxonomy}. However, catalogue
systems usually assume that metadata are supplied by publishers or cataloguers.
They do not by themselves solve the extraction of structured information from
long, heterogeneous web pages, reports, and action sheets.

A second family consists of search and information-retrieval systems. Lexical
engines based on inverted indexes, TF--IDF, or BM25 can efficiently retrieve
documents that contain query terms \citep{salton1988term,robertson2009probabilistic}.
They are well suited to transparent keyword search and can be combined with
facets based on manually supplied metadata. Dense retrieval extends this
approach by representing queries and documents in a shared semantic vector
space \citep{karpukhin2020dense}. Such systems are useful when users express a
need with words that differ from those used in the source documents. Their
limitation for the present project is that retrieval alone does not create the
territorial, organisational, temporal, and thematic metadata required for
comparative analysis.

A third family is constituted by rule-based information-extraction pipelines.
These pipelines use document templates, headings, regular expressions,
lexicons, and pattern-matching rules to identify candidate values. They are
particularly appropriate when source structures are stable because their
decisions can be inspected and corrected. Their limitations have been widely
documented: rules require substantial maintenance and generalise poorly when
documents come from sources with different editorial structures
\citep{sarawagi2008information,grishman2019information}. This limitation is
central in the present corpus, which contains HTML pages, PDF-derived text, and
internally produced action sheets.

More recent systems use transformer-based language models to classify,
summarise, or extract information from text. Prompt-guided LLM extraction makes
it possible to request a predefined JSON schema directly from a document, while
the prompt can specify controlled values and rules for unknown information.
This approach benefits from the contextual modelling capabilities of
transformers \citep{vaswani2017attention,devlin2019bert} and from the flexible
task formulation enabled by large language models \citep{brown2020language}.
It also introduces important risks: a generated field can be fluent and
complete while remaining unsupported by the source. Prompting strategies,
output constraints, and explicit validation are therefore necessary
\citep{liu2023pre}.

The project is not intended to replace a mature public search platform or a
digital library. Rather, it occupies an intermediate position between a
collection pipeline and a future exploration system. Its immediate contribution
is to transform dispersed documentary material into traceable JSON records that
can subsequently support catalogue-like filtering, lexical search, descriptive
statistics, and, in later work, semantic or hybrid retrieval.

\begin{table}[H]
\centering
\begin{tabular}{p{3.2cm}p{3.5cm}p{3.5cm}p{3.4cm}}
\toprule
Approach & Main strength & Main limitation for this corpus & Position of Index\_Insight \\
\midrule
Documentary catalogue & Stable metadata and faceted navigation & Assumes metadata already exist or are manually curated & Reuses the principle of controlled metadata, but automates part of their production \\

Lexical or dense retrieval & Efficient retrieval and exploration of large collections & Does not itself extract reliable documentary metadata & Future use of the structured corpus; not yet a deployed search engine \\

Rule-based extraction & Transparent and inspectable decisions & Fragile across heterogeneous document structures & Baseline and quality-control layer of the pipeline \\

LLM-assisted extraction & Context-sensitive structuring of varied documents & Risk of unsupported inference and output inconsistency & Constrained local extraction, evaluated against human reference annotations \\
\bottomrule
\end{tabular}
\caption{Positioning of Index\_Insight relative to comparable system families.}
\label{tab:comparable-systems}
\end{table}

\section{Synthesis of the state of the art}

The literature shows that no single technique is sufficient for organising a
heterogeneous documentary corpus. Classical information retrieval provides
efficient indexing and transparent lexical search, while dense retrieval can
later improve access to semantically related initiatives. Controlled
vocabularies provide the consistency required for filtering and comparison.
Rule-based extraction remains valuable for traceability and for inspecting
intermediate candidate values, but it becomes difficult to maintain when the
same information is expressed through different layouts and writing styles.
LLMs can reduce this dependence on source-specific rules by interpreting
context, provided that their outputs are constrained and independently checked.

These findings motivate the hybrid design adopted in this thesis. The pipeline
preserves raw sources and intermediate rule-based outputs, then uses a local
LLM to create structured records under a predefined JSON schema and controlled
metadata values. The local deployment supports reproducibility, cost control,
and data governance, while the retained provenance makes it possible to trace a
metadata value back to its source document. Crucially, the system does not treat
model output as self-validating. Completion rates measure technical coverage;
human reference annotation is required to assess factual correctness.

The contribution of \textit{Index\_Insight} is therefore not a claim that an
LLM alone can objectively describe territorial initiatives. It is the design of
a documented workflow that makes a dispersed corpus searchable and comparable,
while preserving the distinction between automatically generated metadata and
validated information. This position also defines the next stage of the
project: after corpus structuring and factual evaluation, the records can be
used to design a faceted, lexical, semantic, or hybrid exploration interface
with the future users of the INSIGHT project.
